\documentclass[10pt,a4paper]{amsart}
\usepackage{amsmath,amssymb,amsthm,graphicx}
% Retained minimal preamble for the Stacks Project chapter more-algebra.tex (More on Algebra).
% The chapter file pinned-original-source.tex carries only \input{preamble} on its first line, so no
% preamble text lives inside the chapter itself.  The Stacks Project uses one single shared preamble for
% every chapter; the snapshot stages it identically as sources/stacks-chow/preamble.tex and
% sources/stacks-etale-cohomology/preamble.tex.  The block below is copied mechanically, line for line,
% from sources/stacks-chow/preamble.tex lines 14-18 (the Xy-pic setup) and 154-217 (the native theorem
% declarations together with the mathematical macros).  Nothing was added, removed or redefined; the
% wrapper class and the amsmath/amssymb/amsthm/graphicx base packages are supplied by the assembler.
% Chapter: more-algebra.tex -- More on Algebra.
% For commutative diagrams we use Xy-pic
\usepackage[all]{xy}

% We use 2cell for 2-commutative diagrams.
\xyoption{2cell}

% Theorem environments.
%
\theoremstyle{plain}
\newtheorem{theorem}[subsection]{Theorem}
\newtheorem{proposition}[subsection]{Proposition}
\newtheorem{lemma}[subsection]{Lemma}

\theoremstyle{definition}
\newtheorem{definition}[subsection]{Definition}
\newtheorem{example}[subsection]{Example}
\newtheorem{exercise}[subsection]{Exercise}
\newtheorem{situation}[subsection]{Situation}

\theoremstyle{remark}
\newtheorem{remark}[subsection]{Remark}
\newtheorem{remarks}[subsection]{Remarks}

\numberwithin{equation}{subsection}

% Macros
%
\def\lim{\mathop{\mathrm{lim}}\nolimits}
\def\colim{\mathop{\mathrm{colim}}\nolimits}
\def\Spec{\mathop{\mathrm{Spec}}}
\def\Hom{\mathop{\mathrm{Hom}}\nolimits}
\def\Ext{\mathop{\mathrm{Ext}}\nolimits}
\def\SheafHom{\mathop{\mathcal{H}\!\mathit{om}}\nolimits}
\def\SheafExt{\mathop{\mathcal{E}\!\mathit{xt}}\nolimits}
\def\Sch{\mathit{Sch}}
\def\Mor{\mathop{\mathrm{Mor}}\nolimits}
\def\Ob{\mathop{\mathrm{Ob}}\nolimits}
\def\Sh{\mathop{\mathit{Sh}}\nolimits}
\def\NL{\mathop{N\!L}\nolimits}
\def\CH{\mathop{\mathrm{CH}}\nolimits}
\def\proetale{{pro\text{-}\acute{e}tale}}
\def\etale{{\acute{e}tale}}
\def\QCoh{\mathit{QCoh}}
\def\Ker{\mathop{\mathrm{Ker}}}
\def\Im{\mathop{\mathrm{Im}}}
\def\Coker{\mathop{\mathrm{Coker}}}
\def\Coim{\mathop{\mathrm{Coim}}}

% Boxtimes
%
\DeclareMathSymbol{\boxtimes}{\mathbin}{AMSa}{"02}

%
% Macros for moduli stacks/spaces
%
\def\QCohstack{\mathcal{QC}\!\mathit{oh}}
\def\Cohstack{\mathcal{C}\!\mathit{oh}}
\def\Spacesstack{\mathcal{S}\!\mathit{paces}}
\def\Quotfunctor{\mathrm{Quot}}
\def\Hilbfunctor{\mathrm{Hilb}}
\def\Curvesstack{\mathcal{C}\!\mathit{urves}}
\def\Polarizedstack{\mathcal{P}\!\mathit{olarized}}
\def\Complexesstack{\mathcal{C}\!\mathit{omplexes}}
% \Pic is the operator that assigns to X its picard group, usage \Pic(X)
% \Picardstack_{X/B} denotes the Picard stack of X over B
% \Picardfunctor_{X/B} denotes the Picard functor of X over B
\def\Pic{\mathop{\mathrm{Pic}}\nolimits}
\def\Picardstack{\mathcal{P}\!\mathit{ic}}
\def\Picardfunctor{\mathrm{Pic}}
\def\Deformationcategory{\mathcal{D}\!\mathit{ef}}
\usepackage{cleveref}
\crefname{originallemma}{original lemma}{original lemmas}
\Crefname{originallemma}{Original lemma}{Original lemmas}
\makeatletter
\newlabel{more-algebra-lemma-helper-finite}{{0EM0}{source}{}{sourcepointer-1}{}}
\newlabel{more-algebra-lemma-helper-finite@cref}{{[originallemma][0][]0EM0}{[1][0][]source}{}{}{}}
\makeatother
\crefname{originallemma}{original lemma}{original lemmas}
\Crefname{originallemma}{Original lemma}{Original lemmas}
\makeatletter
\newlabel{lemma-lift-idempotent-upstairs}{{07M4}{source}{}{sourcepointer-2}{}}
\newlabel{lemma-lift-idempotent-upstairs@cref}{{[originallemma][0][]07M4}{[1][0][]source}{}{}{}}
\makeatother
\crefname{originallemma}{original lemma}{original lemmas}
\Crefname{originallemma}{Original lemma}{Original lemmas}
\makeatletter
\newlabel{algebra-lemma-NAK}{{00DV}{source}{}{sourcepointer-3}{}}
\newlabel{algebra-lemma-NAK@cref}{{[originallemma][0][]00DV}{[1][0][]source}{}{}{}}
\makeatother
\crefname{originallemma}{original lemma}{original lemmas}
\Crefname{originallemma}{Original lemma}{Original lemmas}
\makeatletter
\newlabel{algebra-lemma-integral-integral-over-ideal}{{00H5}{source}{}{sourcepointer-4}{}}
\newlabel{algebra-lemma-integral-integral-over-ideal@cref}{{[originallemma][0][]00H5}{[1][0][]source}{}{}{}}
\makeatother
\crefname{originallemma}{original lemma}{original lemmas}
\Crefname{originallemma}{Original lemma}{Original lemmas}
\makeatletter
\newlabel{algebra-lemma-faithfully-flat-universally-injective}{{05CK}{source}{}{sourcepointer-5}{}}
\newlabel{algebra-lemma-faithfully-flat-universally-injective@cref}{{[originallemma][0][]05CK}{[1][0][]source}{}{}{}}
\makeatother
\crefname{originallemma}{original lemma}{original lemmas}
\Crefname{originallemma}{Original lemma}{Original lemmas}
\makeatletter
\newlabel{algebra-lemma-etale}{{00U2}{source}{}{sourcepointer-6}{}}
\newlabel{algebra-lemma-etale@cref}{{[originallemma][0][]00U2}{[1][0][]source}{}{}{}}
\makeatother
\crefname{originallemma}{original lemma}{original lemmas}
\Crefname{originallemma}{Original lemma}{Original lemmas}
\makeatletter
\newlabel{algebra-proposition-fppf-open}{{00I1}{source}{}{sourcepointer-7}{}}
\newlabel{algebra-proposition-fppf-open@cref}{{[originallemma][0][]00I1}{[1][0][]source}{}{}{}}
\makeatother
\title{A finite algebra over a henselian pair with finite closed fibre}
\author{The Stacks Project Authors}
\date{stacks-project:more-algebra.tex:local-snapshot-20261012; GNU FDL 1.2 or later}
\begin{document}
\maketitle
\setcounter{subsection}{1}
\noindent\textit{Editorial note.} All delivered mathematical text is the author's own, copied from the pinned Stacks Project chapter file more-algebra.tex; no statement or proof step was rewritten, reordered or removed.\par\medskip
\noindent\textit{Editorial note.} The chapter file pinned-original-source.tex carries only \textbackslash{}input\{preamble\} on its first line, so no preamble text lives inside the chapter. The retained file sources/stacks-more-algebra/retained-preamble.tex is a mechanical, line-for-line copy of the project's single shared preamble as staged in the snapshot (Xy-pic setup plus theorem declarations and mathematical macros). Nothing was added, removed or redefined, and the excerpt serves only as preamble, never as a span coordinate.\par\medskip
\noindent\textit{Editorial note.} Both counted spans start at the \textbackslash{}begin line of their environment, so the statement and the author's proof are delivered whole; surrounding prose of the section is deliberately outside the counted spans and is carried as standard prerequisites instead.\par\medskip
\noindent\textit{Editorial note.} The citation fields (source\_id, source\_version, authors, source\_url, license.name, source\_title) are identical to those of every other card this lane has produced from this chapter, so the batch selector treats them as one source.\par\medskip
\noindent\textit{Allowed standard prerequisites (editorial).}\begin{itemize}
\item Henselian pair and the lifting of idempotents. the statement counted by the sibling card S1-algebra-lift-idempotent at 1784-1803. Supplies the splitting of the finite algebra.
\item Finiteness of an algebra over a ring and its preservation under passing to a factor. used at proof lines 2302-2330. Keeps the pieces finite over the pair.
\item Finiteness of the closed fibre and its behaviour under the decomposition. used at proof lines 2330-2344. Supplies the condition that the pieces have to inherit.
\end{itemize}\medskip
\ifdefined\hypertarget\hypertarget{sourcepointer-1}{}\fi
\noindent\textit{Source reference.} Original lemma 0EM0: focused-seeds/AGP/sources/stacks-snapshot-20261011/tags.txt, source lines 18290--18290. The lemma's own canonical tag, quoted so that readers can resolve the label inside the delivered text; it points at the counted statement itself and carries no proof content.\par\medskip
\ifdefined\hypertarget\hypertarget{sourcepointer-2}{}\fi
\noindent\textit{Source reference.} Original lemma 07M4: focused-seeds/AGP/sources/stacks-snapshot-20261011/tags.txt, source lines 9714--9714. Navigation pointer to a statement the counted proof refers to but does not contain; nothing load-bearing is delegated to it.\par\medskip
\ifdefined\hypertarget\hypertarget{sourcepointer-3}{}\fi
\noindent\textit{Source reference.} Original lemma 00DV: focused-seeds/AGP/sources/stacks-snapshot-20261011/tags.txt, source lines 512--512. Navigation pointer to a statement the counted proof refers to but does not contain; nothing load-bearing is delegated to it.\par\medskip
\ifdefined\hypertarget\hypertarget{sourcepointer-4}{}\fi
\noindent\textit{Source reference.} Original lemma 00H5: focused-seeds/AGP/sources/stacks-snapshot-20261011/tags.txt, source lines 631--631. Navigation pointer to a statement the counted proof refers to but does not contain; nothing load-bearing is delegated to it.\par\medskip
\ifdefined\hypertarget\hypertarget{sourcepointer-5}{}\fi
\noindent\textit{Source reference.} Original lemma 05CK: focused-seeds/AGP/sources/stacks-snapshot-20261011/tags.txt, source lines 6923--6923. Navigation pointer to a statement the counted proof refers to but does not contain; nothing load-bearing is delegated to it.\par\medskip
\ifdefined\hypertarget\hypertarget{sourcepointer-6}{}\fi
\noindent\textit{Source reference.} Original lemma 00U2: focused-seeds/AGP/sources/stacks-snapshot-20261011/tags.txt, source lines 1096--1096. Navigation pointer to a statement the counted proof refers to but does not contain; nothing load-bearing is delegated to it.\par\medskip
\ifdefined\hypertarget\hypertarget{sourcepointer-7}{}\fi
\noindent\textit{Source reference.} Original lemma 00I1: focused-seeds/AGP/sources/stacks-snapshot-20261011/tags.txt, source lines 663--663. Navigation pointer to a statement the counted proof refers to but does not contain; nothing load-bearing is delegated to it.\par\medskip
\begin{lemma}
\label{lemma-helper-finite}
Let $(A, I)$ be a Zariski pair. Let $A \to B$ be a finite ring map.
Assume
\begin{enumerate}
\item $B/IB = B_1 \times B_2$ is a product of $A/I$-algebras
\item $A/I \to B_1/IB_1$ is surjective,
\item $b \in B$ maps to $(1, 0)$ in the product.
\end{enumerate}
Then there exists a monic $f \in A[x]$ with $f(b) = 0$ and
$f \bmod I = (x - 1)x^d$ for some $d \geq 1$.
\end{lemma}

\begin{proof}
By Lemma \ref{lemma-lift-idempotent-upstairs} we can find an
\'etale ring map $A \to A'$ inducing an isomorphism $A/I \to A'/IA'$
such that $B' = B \otimes_A A'$ contains
an idempotent $e'$ lifting the image of $b$ in $B'/IB'$.
Consider the corresponding $A'$-algebra decomposition
$$
B' = B'_1 \times B'_2
$$
which is compatible with the one given in the lemma upon
reduction modulo $I$.
The map $A' \to B'_1$ is surjective modulo $IA'$. By Nakayama's lemma
(Algebra, Lemma \ref{algebra-lemma-NAK})
we can find $i \in IA'$ such that after replacing $A'$
by $A'_{1 + i}$ the map $A' \to B'_1$ is surjective.
Observe that the image $b'_1 \in B'_1$ of $b$ 
satisfies $b'_1 - 1 \in IB'_1$.
Thus we may pick $a' \in IA'$ mapping to $b'_1 - 1$.
On the other hand, the image $b'_2 \in B'_2$ of $b$ is in $IB'_2$. By
Algebra, Lemma \ref{algebra-lemma-integral-integral-over-ideal}
there exist a monic polynomial
$g(x) = x^d + \sum a'_j x^j$ of degree $d$ with $a'_j \in IA'$ such that
$g(b'_2) = 0$ in $B'_2$. Thus the image $b' = (b'_1, b'_2) \in B'$
of $b$ is a root of the polynomial $(x - 1 - a')g(x)$. We conclude that
$$
(b' - 1)(b')^d \in \sum\nolimits_{j = 0, \ldots, d} IA' \cdot (b')^j
$$
We claim that this implies
$$
(b - 1)b^d \in \sum\nolimits_{j = 0, \ldots, d} I \cdot b^j
$$
in $B$. For this it is enough to see that the ring map
$A \to A'$ is faithfully flat, because the condition is that
the image of $(b - 1)b^d$ is zero in $B/\sum_{j = 0, \ldots, d} Ib^j$
(use Algebra, Lemma \ref{algebra-lemma-faithfully-flat-universally-injective}).
The map $A \to A'$ flat because it is \'etale
(Algebra, Lemma \ref{algebra-lemma-etale}).
On the other hand, the induced map on spectra is open
(see Algebra, Proposition \ref{algebra-proposition-fppf-open} and use
previous lemma referenced) and the image
contains $V(I)$. Since $I$ is contained in the Jacobson radical of $A$
we conclude.
\end{proof}

\end{document}
